\documentclass[11pt,a4paper]{article}
\usepackage[margin=22mm]{geometry}
\usepackage{amsmath,amssymb,amsthm,booktabs,longtable,array,hyperref,xr}
\externaldocument[U-]{local_actions}
\newtheorem{theorem}{Theorem}[section]
\newtheorem{proposition}[theorem]{Proposition}
\newtheorem{lemma}[theorem]{Lemma}
\newcommand{\F}{\mathbb F}
\newcommand{\U}[1]{\cite[\S\ref{U-#1}]{U}}
\title{Local primitivity in finite flag-transitive linear spaces}
\author{Human Compression project: private unreviewed edition}
\date{}
\begin{document}\maketitle
\begin{abstract}
We give a corrected classification of the faithful flag-transitive
linear-space actions considered by Buekenhout, Dehaye and Leemans.
The large classifications used by that paper remain named external inputs.
The accompanying common paper proves the local-action, field, orbit and
semilinear descent arguments; this application supplies the complete
case assembly, the exact acting-group converses and the finite Feit
bound. Several literal source assertions require corrections. Thus the
result is a clausewise adjudication of the complete source contract,
not a claim that every printed theorem follows unchanged.
\end{abstract}

\section{Objects, source and retained classification inputs}
A linear space here is finite, has a unique line through two distinct
points, at least two points per line, and at least two lines per point.
The group $G\leq\operatorname{Aut}(\Gamma)$ is faithful and
flag-transitive; it need not be the full automorphism group. Write $v,k,r$
for point number, line size and pencil size. Flag transitivity makes $k,r$
constant and $r=(v-1)/(k-1)$. We call the action locally primitive when
$G_x$ is primitive on its $r$ lines and $G_L$ is primitive on its $k$
points, taking the actual induced actions and their kernels. For the
source's definition of residual weak primitivity (RWPRI), this is
equivalent by the retained Higman--McLaughlin point-primitivity theorem;
the primitive-stabilizer criterion is the standard Dixon--Mortimer fact.
Local two-transitivity means two-transitivity of both these actions and
is the source's $(2T)_1$ condition.

The primary edition is the complete twenty-two-page article
\cite{BDL}, printed pages 25--46. Its bibliography and theorems quoted
from earlier authors are part of its dependency boundary. We retain the
following inputs, at their actual permutation-representation scope.
Their large classification proofs are not claimed as new proofs here.

\begin{enumerate}
\item BDDKLS, with the subsequent classification work explicitly listed
by the source (Liebeck, Delandtsheer, Kleidman, Saxl and the sporadic
cases), reduces a thick non-one-dimensional flag-transitive linear
space to a projective space, affine space, one of the two Hering spaces,
a Witt--Bose--Shrikhande space, a Hermitian unital or a Ree unital
\cite[Proposition 1.1 and its attribution]{BDL}.
\item Kantor and Cameron--Kantor's projective classification gives the
classical special-linear case, the specified $A_7$ action on
$\operatorname{PG}(3,2)$, or the sharply flag-transitive Frobenius
plane. Feit's exact arithmetic restrictions on the last case are
retained \cite[Propositions 3.1--3.2]{BDL}.
\item For affine dimension at least three, retain the precise
BDDKLS--Huybrechts alternatives and their natural field-module
identifications: the one-dimensional case; the impossible
last-derived-$2.A_5$ case on $\operatorname{AG}(4,3)$; or a
point-two-transitive action with one of the six representation types
in source Proposition 3.4. The normal natural groups there are
$\operatorname{SL}_u(q^{n/u})$, $\operatorname{Sp}_u(q^{n/u})$
or $G_2(q^{n/6})'$, together with the three listed small exceptions.
Their invariant extension-field structures are part of these
representation identifications, not an inference from an abstract
isomorphism type.
\item Retain Foulser's actual affine-plane representations. These include
the ten groups on $\operatorname{AG}(2,9)$ with common minimal
subgroup $3^4:2.A_5$. The induced special-linear group is the projective special-linear group,
not necessarily the full projective linear group. The two printed
inclusions are corrected below.
\item The BDD non-one-dimensional non-Desarguesian planes are the
L\"uneburg planes, the nearfield plane of order nine and the Hering
plane of order twenty-seven. Retain their specified full-group failures:
L\"uneburg is not locally point-primitive; the latter two are not
locally line-primitive. Retain the two $3^6$-point Hering spaces of
line size nine with group $3^6:\operatorname{SL}_2(13)$, locally
primitive but not locally two-transitive; and the Ree-unital full-group
failure of local line-primitivity \cite[\S\S3.2.2, 3.3, 3.6]{BDL}.
\item Retain the full automorphism identifications for the WBS and
Hermitian geometries, respectively $\operatorname{P\Gamma L}_2(2^e)$
and $\operatorname{P\Gamma U}_3(q)$, with the classical natural
models. Their required local actions are proved in the common paper,
including the small-parameter exceptions to the source assertions.
\item For circles, retain the complete almost-simple and affine
point-two-transitive representation roster quoted as source
Proposition 4.1, with the original exact Foulser exception list. Also
retain the specified Mathieu transitivity, Atlas stabilizer/action
identifications for the sporadic and small alternating cases, and the
classical orthogonal-point symplectic subdegrees. These local
representation inputs are stated explicitly in Section~\ref{circles}.
\item Delandtsheer's one-dimensional reduction supplies all translations
and the spread/nonspread dichotomy. Retain the full unordered
intersecting-line-pair classification in source Proposition 5.1(ii),
the nonspread divisibility in Proposition 5.7 and the more general
prime-power repunit theorem quoted from Ribenboim in Proposition 5.3.
Our own prime-pencil proof uses only the elementary prime case and
is written in full; it does not reprove the quoted prime-power theorem.
For precision, the retained unordered-pair exceptions are: the circles
with $\operatorname{AGL}_1(3)$, $\operatorname{AGL}_1(4)$,
$A\Gamma L_1(4)$ or $A\Gamma L_1(8)$; the Fano plane with $7:3$;
and $\operatorname{AG}(2,4)$ with $A\Gamma L_1(16)$.
The repunit statement retained is that
$(x^u-1)/(x-1)=r^m$, $u\geq3$, $x>1$, $r$ prime, forces $u$
prime and the multiplicative order of $x$ modulo $r$ equal to $u$;
if $x=s^b$ as in the quoted statement, then $b=u^e$ and
$r\equiv1\pmod{u^{e+1}}$, with $r$ not a Fermat prime.
Its prime specialization below is separately proved.
\end{enumerate}
The natural $G_2$ module and maximal parabolic identification are retained
classical building inputs. Its binary derived-group boundary uses one
complete small GAP maximal-subgroup calculation. The large classification
inputs, missing original software and nontext graphics have no asserted
zero cost or full formalization.

\section{Thick actions and the affine comparison}
The common local-action criterion identifies a translation pencil with
its spread, and the line action with $U\rtimes H_U|_U$. Thus the two
conditions are precisely pencil primitivity and irreducibility of the
\emph{actual restriction} $H_U|_U$; for local two-transitivity they
are pencil two-transitivity and transitivity on $U\setminus\{0\}$
\cite[Proposition~\ref{U-translation}]{U}.
The distinction matters when a line stabilizer has a restriction kernel.

Classical projective and affine special-linear actions are locally
two-transitive by the explicit basis and determinant arguments in
\cite[Proposition~\ref{U-classical}]{U}. In the exceptional projective
$A_7$ action the retained local groups are $\operatorname{PSL}_3(2)$
on seven pencil points and $S_3$ on a line; both are two-transitive.
The other projective branch reduces, by the written Feit congruence
argument, to the Fano plane with $7:3$; its cyclic local groups of
order three are primitive but not two-transitive
\cite[Proposition~\ref{U-projectivefrobenius}]{U}.

For the affine classification one cannot assume that two geometries
are nested merely because their line sizes are ordered. The full
replacement for source Lemma 3.5 is
\cite[Theorem~\ref{U-affinemax}]{U}: a primitive affine pencil has
largest line size among all invariant affine spaces for the same
permutation group. The proof treats non-Desarguesian planes as well:
the $\operatorname{SL}$ and $\operatorname{Sp}$ vector stabilizers
have outside orbits larger than the affine bound $\sqrt v$; the
$G_2$ case uses its calculated commuting field and the named BDD
L\"uneburg exclusion. For an alternative line through two fixed
points, uniqueness makes it invariant under their stabilizer. The
argument neither postulates a common field in advance nor identifies
$G$ with a full automorphism group.

More directly, extension-field normalizers preserve the actual coarser
field-line spread. The commuting-algebra lemma therefore excludes all
proper-field natural modules from a primitive pencil. The remaining
natural symplectic pencil is primitive and has two nontrivial
perpendicularity orbits. Its whole normalizer preserves that relation,
so no outer extension makes it two-transitive
\cite[Proposition~\ref{U-symplectic}]{U}. The same normalizer argument
and the retained maximal-parabolic action treat $G_2(q)$, $q\geq4$
even \cite[Proposition~\ref{U-naturalg2}]{U}. For $q=2$ use
\cite[Proposition~\ref{U-binaryg2}]{U}, with the required complete
order list $6048;96,96,168,216$ and automorphism order $12096$.
Its degree-63 pencil is primitive and not two-transitive.

The other Huybrechts cases have imprimitive pencils by the explicit
normal-subgroup/order exclusions in the common paper. The Foulser
plane exceptions reduce to $q=9$: a faithful primitive $A_4$ or
$S_4$ action has degree at most four, and the primitive degrees of
$A_5$ are $5,6,10$. The retained ten degree-81 groups are all locally
primitive. In their minimal $2.A_5$ linear group an order-four lift
$f$ of an involution has $f^2=-I$ and fixes a vector line over $\F_9$;
its order-four scalar acts irreducibly on that two-dimensional
$\F_3$-line. This supplies the line-primitivity proof without the
source's unjustified assertion that the whole Sylow subgroup fixes the
same directions \cite[Proposition~\ref{U-nineplane}]{U}.
Their pencil images $A_5,S_5$ cannot be two-transitive of degree ten
by order divisibility.

The retained non-Desarguesian-plane and Ree exclusions pass from full
groups to every subgroup: a primitive subgroup would force primitivity
of the full local action. The Hering spaces retain precisely their
stated BDD local properties. For WBS spaces the uniform common proof
is stronger than a finite check but corrects the source's stated reason:
if $q+1$ is composite the pencil is imprimitive; if it is prime,
$e$ is a power of two and the line action on $\F_{2^e}/\F_2$
has reducible cyclic two-group stabilizer. Thus no WBS space is
locally primitive. At $q=8$ the line action of the full group is
$A_4$ of degree four and \emph{is} primitive; its nine-point pencil
fails instead \cite[Propositions~\ref{U-wbs-model}--\ref{U-wbs-exclusion}]{U}.

For the Hermitian unital, put $d=\gcd(3,q+1)$ and
$\mathcal C=\F_{q^2}^*/(\F_{q^2}^*)^d$. The common matrix
calculation defines the exact outer action
$\pi:G/\operatorname{PSU}_3(q)\to\operatorname{Sym}(\mathcal C)$.
The pencil group in the special-unitary case is
$\F_{q^2}^+\rtimes(\F_{q^2}^*)^d$. For $q>2$ it is primitive;
for $q=2$ it is not. In either case the full acting group is locally
two-transitive exactly when $\pi(G)$ is transitive. For $q=2$ this
same condition is necessary for local primitivity. All line actions
are already two-transitive in the special-unitary subgroup
\cite[Theorem~\ref{U-unitary}]{U}. In particular the literal blanket
converse for every $\operatorname{PSU}_3(q)$-containing group is false.

\section{The complete corrected circle test}\label{circles}
For a circle, flags are ordered pairs of distinct points. Thus flag
transitivity is point two-transitivity, RWPRI is primitivity of the
point stabilizer on the remaining points, and $(2T)_1$ is point
three-transitivity \cite[Proposition~\ref{U-circles}]{U}. The following
list is necessary and sufficient within the retained two-transitive
representation classification. Repeated or overlapping small actions
are permitted; the entries describe actions, not just group names.

\begin{longtable}{@{}p{.60\textwidth}p{.34\textwidth}@{}}
\toprule
Locally primitive circle action & Locally two-transitive exactly when\\\midrule\endhead
$A_v$ or $S_v$, $v\geq5$ & Always.\\[2pt]
$\operatorname{PSL}_2(q)\leq G\leq\operatorname{P\Gamma L}_2(q)$,
$v=q+1$ & $q$ even, or an allowed semilinear representative has
nonsquare determinant.\\[2pt]
$\operatorname{Sp}_{2n}(2)$, $v=2^{2n-1}\pm2^{n-1}$, $n\geq3$
& Never.\\[2pt]
$\operatorname{PSL}_2(11)$, $v=11$ & Never.\\[2pt]
$M_v$, $v=11,12,22,23,24$; $\operatorname{Aut}(M_{22})$, $v=22$;
$M_{11}$ in degree twelve & Always.\\[2pt]
$HS$ in degree $176$; $Co_3$ in degree $276$ & Never.\\[2pt]
$2^n:\operatorname{GL}_n(2)$, $n\geq2$ & Always.\\[2pt]
$2^n:H$, $n\geq4$ even, with natural
$\operatorname{Sp}_n(2)\leq H\leq N_{\operatorname{GL}_n(2)}
(\operatorname{Sp}_n(2))$ & Never.\\[2pt]
$2^4:A_6$ and $2^4:A_7$ in their retained natural binary actions
& Exactly the $A_7$ case.\\[2pt]
$2^6:H$, $G_2(2)'\leq H\leq N_{\operatorname{GL}_6(2)}(G_2(2)')$
& Never.\\[2pt]
One-dimensional, point-two-transitive $G\leq A\Gamma L_1(v)$,
with $v-1$ prime & Exactly $v=3,G=S_3$, or $v=4,G=S_4$.\\
\bottomrule
\end{longtable}

Here are the complete local justifications, including the retained
small-action boundaries. Alternating groups in the stated range are
three-transitive by even-permutation completion of a prescribed triple;
for $A_5$ the two unused letters already suffice. The projective-line
entry follows from the scalar-square quotient calculation in
\cite[Proposition~\ref{U-projectiveline}]{U}; it includes mixed
field-diagonal extensions and does not require containment of
$\operatorname{PGL}_2(q)$ when that would be too strong.
The source's formulation via normalization of a sharply three-transitive
group is retained as the attributed Gardiner--Praeger--Zhou description
of these same extensions; our determinant criterion tests the actual
specified action.

For the symplectic orthogonal-point actions, retain the rank-three
subdegrees and use the complete block-size divisibility proof in
\cite[Proposition~\ref{U-orthogonalcircles}]{U}. For
$\operatorname{PSL}_2(11)$ the retained point stabilizer is $A_5$ on
ten points, with stabilizer $S_3$. The elementary maximality proof in
\cite[Proposition~\ref{U-nineplane}]{U} gives primitivity, while
$10\cdot9\nmid60$ rules out two-transitivity. The stated Mathieu
transitivity is the retained exact permutation-action input. In the
$HS$ action the retained Atlas point stabilizer is
$\operatorname{PSU}_3(5):2$, primitive but not two-transitive on the
other175 points. In the $Co_3$ action the retained two-point inclusion
$\operatorname{PSU}_4(3):2<McL:2$ is maximal and its degree-275
action is not two-transitive. These are explicit retained local-action
claims, not order-only inferences.

For binary affine groups, translations reduce the question to the
nonzero-vector action of the point stabilizer. The natural linear and
symplectic cases have already been proved in the common paper. The
retained $A_6\leq\operatorname{GL}_4(2)$ action has point stabilizer
$S_4$ maximal of index fifteen; it is primitive, and
$15\cdot14\nmid360$ excludes two-transitivity. The retained
$A_7\leq\operatorname{GL}_4(2)$ action is the two-transitive action
on the fifteen points of $\operatorname{PG}(3,2)$, giving
three-transitivity of its affine extension. This does not make the
\emph{fifteen-point circle} of $A_7$ primitive: there the point
stabilizer preserves seven pairs, the punctured projective lines.
The binary $G_2(2)'$ case is the exact small case proved above.

In the one-dimensional entry, the point stabilizer is cyclic by cyclic.
Its primitive remaining-point action therefore has prime degree
$r=v-1$, and this condition is sufficient by prime-degree transitivity.
Two-transitivity on those $r$ points requires
$r(r-1)\mid r d$ for $v=p^d$, so $p^d-2\mid d$.
For $d\geq3$, $2^d-2>d$; for $d=2$ only $p=2$ works, and for
$d=1$, $v\geq3$ leaves only $p=3$. The required group orders force
$S_4$ and $S_3$, respectively, and both suffice.

For exhaustion, run through the retained roster in source
Proposition~4.1. Higher-dimensional projective and unitary actions
preserve actual thick linear spaces; their punctured lines form
nontrivial invariant pencil partitions. Suzuki and Ree point stabilizers
have regular unipotent groups with proper nontrivial characteristic
centers, giving local block systems even in the outer extensions.
Odd-characteristic affine modules of dimension at least two preserve
thick prime-field affine lines. For even extension fields larger than
two the actual normalized extension-field lines give the same
exclusion. The remaining extraspecial and $\operatorname{SL}_2(5)$
or $\operatorname{SL}_2(13)$ affine exceptions have odd characteristic
and are excluded this way. Over the binary field the cases in the
table remain. This accounts for every branch of the retained roster.
The small circles are included exactly: degree three requires $S_3$;
degree four permits $A_4,S_4$ for local primitivity and only $S_4$
for local two-transitivity. Literal extensions of the source's
alternating-group footnotes to $A_3$ or to a three-transitive $A_4$
are invalid.

\section{One-dimensional spaces and the acting-group converse}
In the spread case, the complete semilinear descent proof in
\cite[Theorem~\ref{U-spread}]{U} gives a Desarguesian spread with
\[
 v=p^{u^{e+1}},\qquad k=p^{u^e},\qquad
 r=(v-1)/(k-1)\ \hbox{prime},\qquad p,u\ \hbox{prime},\quad e\geq0.
\]
It proves the earlier source question affirmatively at this precise
one-dimensional, flag-transitive, locally primitive scope. That answer
is reused from the accepted phase-four proof, with no new question credit.
For a specified acting group $G=F^+\rtimes H\leq A\Gamma L_1(F)$,
flag transitivity and the prime pencil do \emph{not} alone suffice.
The exact further condition is that $H_K|_K$ be irreducible over
$\F_p$, for the line $K=\F_{p^{u^e}}$ through zero. This is both
necessary and sufficient by the translation criterion. Existence for
every allowed parameter set is proved using the actual one-dimensional
Singer affine group $F^+\rtimes F^*$, rather than borrowing the full
higher-dimensional automorphism group.

An explicit counterexample to the source's unrestricted acting-group
converse is $F=\F_{16}$, $K=\F_4$ and
$G=F^+\rtimes C_5$, with $C_5\leq F^*$.
The five $K$-directions form a regular $C_5$ orbit, so $G$ is
flag-transitive on $\operatorname{AG}(2,4)$ and satisfies the
parameters of source Theorem~6.1(1). But
$C_5\cap K^*=1$, so its action on a line is just the regular
additive group $K^+\cong C_2^2$. It preserves the two cosets of any
order-two subgroup and is imprimitive. This corrects the printed
converse, without affecting the existential assertion of Proposition~5.5.

The local-two-transitive spread cases are exactly
$\operatorname{AG}(2,2)$ with $A\Gamma L_1(4)$ and
$\operatorname{AG}(2,4)$ with $A\Gamma L_1(16)$; all necessity,
uniqueness of the acting group and sufficiency proofs are in
\cite[Proposition~\ref{U-onedim2t}]{U}. The nonspread proof
\cite[Proposition~\ref{U-nonspread}]{U} leaves exactly the Fano
plane with $7:3$ and the three-point circle with $S_3$; only the
latter is locally two-transitive. The source's Proposition~5.8 and
its footnote removing the spread restriction miss that small circle
when passing from $r=k$ to a projective plane. Both final source
unions already contain it through the degree-three projective-line
circle, so no additional family beyond their overall universe is needed.

\section{The final classification with both directions}
\begin{theorem}[Corrected locally primitive classification]\label{finalprime}
Subject to the named retained classifications above, the faithful
flag-transitive locally primitive finite linear-space actions are
exactly the following union, with their specified natural representations.
\begin{enumerate}
\item The one-dimensional Desarguesian spread actions of Section~4,
with prime pencil and the actual line-stabilizer irreducibility condition.
\item $\operatorname{PG}(n,q)$, $n\geq2$,
$\operatorname{PSL}_{n+1}(q)\leq G\leq\operatorname{P\Gamma L}_{n+1}(q)$;
the specified $A_7$ action on $\operatorname{PG}(3,2)$; and the Fano
plane with $7:3$.
\item $\operatorname{AG}(n,q)$, $q\geq3$, with translations and
normal natural $\operatorname{SL}_n(q)$, $n\geq2$, or
$\operatorname{Sp}_n(q)$, $n\geq4$ even, in its linear stabilizer.
\item $\operatorname{AG}(6,q)$, $q$ even, with translations and
normal natural $G_2(q)'$ in its linear stabilizer, including $q=2$.
\item The ten specified Foulser groups on $\operatorname{AG}(2,9)$
with pencil image $A_5$ or $S_5$.
\item The two retained Hering spaces on $3^6$ points with line size
nine and group $3^6:\operatorname{SL}_2(13)$.
\item A Hermitian unital, with
$\operatorname{PSU}_3(q)\leq G\leq\operatorname{P\Gamma U}_3(q)$:
all such actions for $q>2$; for $q=2$ precisely those whose scalar-coset
action $\pi(G)$ is transitive.
\item The circle actions in the complete table of Section~\ref{circles}.
\end{enumerate}
All stabilizers in the affine entries are actual semilinear geometry
stabilizers normalizing the stated natural subgroup. The entries can
overlap; no claim of pairwise disjointness is made.
\end{theorem}
\begin{proof}
For a thick non-one-dimensional space apply the retained six-family
BDDKLS reduction. The projective local calculation gives entry~2.
The full affine comparison and the explicit pencil exclusions leave
entries~3--5; the non-Desarguesian-plane cases fail one of the retained
full-group local tests. The retained Hering local actions give entry~6.
The uniform WBS proof and retained Ree exclusion remove those two
families. The unitary scalar-coset calculation gives entry~7, with
its necessary small-parameter qualification. The one-dimensional
spread and nonspread arguments give entry~1 and the already listed
Fano/three-point cases. Circles give exactly entry~8 by the exhaustion
proof in Section~\ref{circles}. These arguments prove necessity.

Conversely, the translation criterion plus the explicitly stated
restriction irreducibility proves entry~1. The basis, Frobenius,
symplectic, exceptional-pencil and order-four-lift proofs in the common
paper give entries~2--5 for every stated group; supergroups preserve
local primitivity. Entry6 uses precisely its named retained BDD
local-action assertion. The unitary calculation gives entry~7 for
exactly the listed groups. Section~\ref{circles} proves every circle
entry. Each action is faithful and flag-transitive in its specified
model; global point-primitivity is supplied by the retained
Higman--McLaughlin theorem or by the direct translation proof.
Thus both directions and every local hypothesis are established.
\end{proof}

\begin{theorem}[Corrected locally two-transitive classification]\label{final2t}
Under the same retained inputs, the locally two-transitive actions are
exactly the following subunion of Theorem~\ref{finalprime}:
$\operatorname{AG}(2,4)$ with $A\Gamma L_1(16)$;
the classical projective actions and the projective $A_7$ action;
the natural affine special-linear actions with $q\geq3$;
the Hermitian actions for which $\pi(G)$ is transitive;
and exactly the circle entries marked locally two-transitive in the table.
\end{theorem}
\begin{proof}
Local two-transitivity implies local primitivity, so start with the
complete preceding union. The one-dimensional proof leaves only the
four- and sixteen-point cases, and the former is already a circle.
The classical projective and affine actions and projective $A_7$
action are locally two-transitive by their explicit local groups.
The Fano cyclic action, symplectic perpendicularity relation, exceptional
$G_2$ relation/order bound, Foulser degree-ten order test and retained
Hering local assertion exclude all other thick entries except unitals.
There the necessary and sufficient condition is exactly the calculated
transitivity of $\pi(G)$. The circle criterion and table finish both
directions. This proves the classification for the specified acting
groups, including the converses that were only asserted in the source.
\end{proof}

\section{The finite Feit bound stated in the source}
The source's physical page~5 reports a computer-assisted extension to
plane order $m\leq10^{10}$ of Feit's exclusion. This is a separate finite
assertion, not a consequence of local primitivity, and we retain it explicitly.
The following calculation recovers the stated bound while using Feit's
necessary conditions as an attributed external theorem.

\begin{proposition}[The declared finite range]\label{feitbound}
Assume the retained Kantor--Cameron--Kantor classification and Feit's theorem
in source Propositions~3.1--3.2. A flag-transitive projective plane of order
$m\leq10^{10}$ is Desarguesian.
\end{proposition}
\begin{proof}
The retained classification reduces a possible non-Desarguesian example
to its sharply flag-transitive Frobenius case. Feit's necessary conditions
then give
\[
 8\mid m,\quad m\text{ not a power of }2,\quad
 v=m^2+m+1\text{ prime},\quad d^{m+1}\equiv1\pmod v
 \quad(d\mid m).
\]
The order of $m$ modulo $v$ is exactly three: $m^3\equiv1$ and
$1<m<v$ rules out order one. The condition with $d=m$ therefore gives
$3\mid m+1$. Combining with $8\mid m$ gives $m=8+24t$. The condition
with $d=2$ gives $2^{m+1}\equiv1\pmod v$.

The complete required program \texttt{feit\_bound.cpp} checks all
$416\,666\,667$ integers of the form $8+24t$ in the declared range.
It generates the primes $p\leq10^4$ by the ordinary Eratosthenes sieve
and enumerates, directly, all residue classes $t\bmod p$ for which
$(8+24t)^2+(8+24t)+1\equiv0\pmod p$. Segmented marking rejects a value
only when that $p$ is a proper divisor of $v$; the small possible equality
$v=p$ is explicitly retained. Thus every marked rejection contradicts the
required primality, with no probabilistic primality test. Powers of two
are separately excluded by the exact bit test. On every remaining value,
GMP computes $v$ and $2^{m+1}\bmod v$ as arbitrary-precision integers.
The complete saved result is
\[
\begin{array}{lr}
\text{proper-divisor rejections}&331\,513\,606\\
\text{excluded powers of two}&16\\
\text{remaining exact modular powers}&85\,153\,045\\\hline
\text{total}&416\,666\,667.
\end{array}
\]
None of the remaining modular powers equals one. Every possible example
therefore violates a necessary condition.

The loop parameters use 64-bit unsigned integers and never exceed
$10^{10}+1$. The only native squaring of an unreduced $m$ occurs behind
the guard $m\leq10^4$ when testing $v=p$. The large polynomial and modular
power are formed in GMP; near the upper boundary the modulus needs 67
bits and is not silently truncated to 64. The full execution receipt
records the exact code hash, commands, types and counts. A separate
109-value check, including both sides of $2^{32}$ and the upper endpoint,
compares this GMP arithmetic with Python's unbounded integers. Direct
unsieved Python testing through $10^6$ gives the same empty necessary
candidate list. These latter checks are optional validation; the full
exact bounded run and the retained Feit theorem are the proof premises.
\end{proof}

This recovers the authors' finite numerical remark, not a new unbounded
answer about non-Desarguesian flag-transitive planes. The original Magma
and GAP source code was not supplied and its complete cost is unknown.
The complete current C++ file and its operational runner will be disclosed
separately under the symmetric executable-code policy.


\section{Literal source corrections and scope}
The source's classification strategy and its retained inputs remain
attributed to their authors. The present common interfaces replace
repeated local arguments and make the acting-group conditions explicit.
For clarity, the complete source map records the following corrections,
in addition to the retained quoted results and recovered proofs.
\begin{enumerate}
\item Lemma~3.5 has a new complete proof at its full affine scope;
its printed automatic-nesting and full-group-equality sentences are
not used. This is a proof repair, not a theorem counterexample.
\item The induced groups in Propositions 3.7--3.8 contain
$\operatorname{PSL}_2(q)$. Containment of $\operatorname{PGL}_2(q)$
is not necessary.
The explicit $\operatorname{ASL}_2(3)$ example distinguishes them.
The Proposition 3.9 conclusion survives: the order-four lift proof
replaces its Sylow fixed-direction shortcut.
\item The WBS full group need not fail local point-primitivity:
$W(8)$ has line action $A_4$. The uniform no-RWPRI conclusion remains
true, with the pencil failure included. Its retained identification
with the smallest Ree unital is unchanged.
\item The Hermitian blanket local-two-transitivity assertion is false
when $\gcd(3,q+1)=3$ for the special-unitary group; at $q=2$ even
local primitivity fails for that group. The exact scalar-coset criterion
replaces both affected converses.
\item The $A_7$ fifteen-point circle must be deleted from the source's
primitive and three-transitive lists. The source's binary $G_2(2)'$
circle exclusion incorrectly treats a two-point line as thick; its
primitive degree-63 pencil survives. One-dimensional primitive circles
omitted from the stand-alone Proposition~4.3 are included explicitly.
In Theorem~6.2 the explicit three-transitivity antecedent makes its
$A_7$ row empty; that redundant necessary row alone is not a
counterexample to the theorem's implication. The binary exceptional
case is already in the source's final affine
$G_2$ branch, so its restored circle description does not enlarge that
final union. Small alternating-group footnotes require the exact
three- and four-point treatment given above.
\item Proposition~5.8 and the footnote after Proposition~5.6 need the
three-point nonspread circle. The full acting-group converse in
Theorem~6.1(1) needs line-restriction irreducibility, as witnessed by
$2^4:C_5$; the geometric uniqueness question is settled only at the
already accepted one-dimensional locally primitive spread scope.
\end{enumerate}
The source's stated $q=2$ unitary/affine overlap is preserved:
$\operatorname{ASL}_2(3)$ and $\operatorname{AGL}_2(3)$ occur on the
nine-point affine plane/unital in both descriptions. This is an overlap
of specified actions. We do not promote the source's tentative remark
about intersections of families to a new exhaustive disjointness claim.

There are two required finite premises: the complete small
$\operatorname{PSU}_3(3)$ maximal-subgroup calculation and the exact
bounded Feit run. The WBS8 computation is optional validation of the
uniform written proof. Operational compilation/render/checking files
are inventoried separately. Source software is mentioned but not
supplied; its full executable cost is unknown. Any later measurement
must apply symmetric executable boundaries and retain that unknown,
not count only our software as a compression factor. No Lean theorem
is added and no end-to-end formalization is claimed. This private
edition has no source-admission or numerical status before distinct review.

\begin{thebibliography}{9}
\bibitem{BDL} F. Buekenhout, P.-O. Dehaye and D. Leemans,
\emph{RWPri and $(2T)_1$ flag-transitive linear spaces},
Beitr\"age zur Algebra und Geometrie \textbf{44} (2003), 25--46.
All external classification locators above refer to this fixed edition's
numbered statements and its complete original bibliography.
\bibitem{U} \emph{Local actions, invariant pencils and scalar quotients},
companion common paper in this private edition. It includes the full
previously reviewed phase-four semilinear descent proof, with that
prior work and its dependencies explicitly attributed.
\end{thebibliography}
\end{document}
